\chapter{Where to go next}

\section{The documentation you already have}

The in-tree corpus is over 900 files and it is the primary source for
everything in this book. From any prompt:

\begin{repl}
: help matches          ;;; practical, task-shaped
: ref  vmcode           ;;; the exhaustive reference -- and the VM spec
: teach json            ;;; tutorials, including how lib json was built
\end{repl}

The same material is published at
\code{https://iotone.github.io/poplog/}, regenerated on every push by a
Pop-11 program; \code{llms.txt} there is for AI assistants. The handful of
entries worth reading early: \code{help external} and \code{ref external}
(Chapter~7), \code{ref vmcode} and \code{ref subsystem} (Chapter~6),
\code{ref regexp}, \code{help json}, \code{help crypto}, and
\code{teach swank}.

\section{Four decades of teaching material}

Poplog was the vehicle for a great deal of university AI teaching, and most
of it survives. The tree does not vendor any of it; one script fetches it
from public archives and wires it into \code{help}/\code{teach} inside the
system:

\begin{sh}
tools/fetch-learning.sh list        # what's available
tools/fetch-learning.sh --all       # about 30 MB into learn/
\end{sh}

\begin{center}
\footnotesize
\begin{tabular}{lp{8.1cm}}
\toprule
\code{hidden-gems} & Small idiomatic demos of what makes Pop-11 distinctive:
  the open stack, updaters, dynamic lists, coroutines, \code{dlocal}, the GC,
  extensible syntax.\\
\code{bhamteach} & The Birmingham introductory-AI course --- pattern-matching
  chatbots, grammars and parsing, story generation, recursion exercises.\\
\code{packages} & The classic package tree: SimAgent and Poprulebase,
  popvision, rclib, neural nets, Braitenberg vehicles.\\
\code{paradigms} & A UMASS programming-language-paradigms course taught in
  Pop-11 (1997--2000).\\
\code{primer} & The Pop-11 Primer --- the standard text for experienced
  programmers.\\
\code{examples} & Curated teaching examples: Othello, Life, robolang.\\
\bottomrule
\end{tabular}
\end{center}

Upstream licences apply per pack; each records its origin in
\code{learn/\emph{pack}/.source}.

\section{Reading the source}

Poplog is written mostly in itself, which makes it unusually approachable.

\begin{itemize}
\item \code{pop/src/} --- the core: \code{vm\_plant.p} and \code{vm\_control.p}
  are the VM builder from Chapter~6; \code{syscomp/} holds one directory per
  CPU back-end.
\item \code{pop/lib/lib/} --- the library, and the best source of idiomatic
  Pop-11 at length. \code{json.p}, \code{crypto.p}, \code{http\_server.p},
  \code{jsonrpc.p}, \code{swank.p} and \code{forth.p} are all readable in an
  afternoon each.
\item \code{pop/plog/}, \code{pop/lisp/}, \code{pop/pml/}, \code{pop/forth/}
  --- the four front-ends, i.e.\ four worked answers to Chapter~6.
\item \code{examples/} --- runnable programs: the HTTP service, ELIZA, a
  Z-machine interpreter, graphics demos.
\end{itemize}

\section{Benchmarks, ports, and the state of the project}

\code{BENCHMARKS.md} carries the full cross-platform, cross-language matrix
with its methodology; the numbers quoted in Chapter~1 are drawn from it. The
\code{PORTING-*.md} files are the working notes from the AArch64, Apple
Silicon and RISC-V ports --- including the bugs, which is the part usually
missing from porting documentation and the part most useful to the next
porter. Windows is the one platform still marked TODO.

\begin{callout}{If you take one idea away}
It is that a compiler cheap enough to call in a loop changes what a
programming environment can be. Everything distinctive here --- live
redefinition, saved images, five languages in one heap, an editor that talks
to a running process, an agent that keeps a session warm across tool calls
--- falls out of fifteen microseconds per definition.
\end{callout}
